\clearpage\fancyhead[L]{\small BORDA’S CIRCLE.}\hypertarget{native-409-part-lower}{}\label{native-409-part-lower}
\begin{center}\large BORDA’S CIRCLE.\end{center}
\begin{center}\hrule\vspace{4pt}\fontsize{10}{12}\selectfont V. \textsc{Liadenau} and \textsc{Bohnenberger}, Zeitschrift für Astronomie u. verwandte Wiss.\\ Vol. V. pp. 198..211. March and April 1818.\par\vspace{4pt}\hrule\vspace{4pt}\end{center}
\par\indent As you requested, my most esteemed Professor, I am pleased to share with you here the method I have used to set the line of sight parallel to the plane of the circle. Repeated experience has taught me that it can achieve very high precision.

\par\indent To express myself concisely and precisely, I designate the different axes of the instrument as follows. I call the axis that must be set vertically for altitude measurements, and that is kept at least fairly close to this position for every other use, the \emph{first axis}. Motion about this axis is measured on the azimuth circle. This circle, however, has only one index, whose vernier gives only minutes, though fractions of a minute can still be estimated. Moreover, it is uncertain whether there is a slight eccentricity. So an adjustment of the line of sight using the azimuth circle (by determining the azimuth of a terrestrial object at two opposite vertical positions of the main circle, allowing for its parallax if necessary when the object is not very far away) cannot yet yield all the precision that is attainable or desirable. This convenient and simple procedure may still be used for a first preliminary adjustment.

\par\indent Nearly perpendicular to the first axis is the one I call the \emph{second}. It is therefore always nearly horizontal. Rotation about it brings the circle into a vertical, inclined or horizontal position. The \textsc{Reichenbach} circles have no graduated scale for measuring the amount of this rotation. Such a scale is unnecessary for most observations, but there are a few cases in which it would be desirable. In particular, angles between fixed stars and terrestrial objects could easily be measured in daylight, whereas at present this is either impossible or possible only with great difficulty.

\par\indent Again, nearly perpendicular to the second axis is the one about which the whole circle turns. This motion is measured on the drum circle. The latter axis coincides, in the mathematical sense, with the one \SourceLeaf{410} about which the alidade circle moves, so the two axes are treated as one, which I call the \emph{third}. This third axis is exactly perpendicular to the plane of the circle. If $a$ is the angle that the line of sight makes with the plane of the circle (which I take as positive when the objective end is farther from the face of the circle than the eyepiece end), then $90^\circ{-}a$ is the angle between the line of sight and the third axis, with the line of sight taken toward the objective end and the third axis taken downward from the graduated face.

\par\indent To use the instrument, the second axis need not be exactly parallel to the plane of the circle or exactly perpendicular to the first axis. The deviations will certainly be small, and repeated experience with the 12-inch \textsc{Reichenbach} circle here has taught me that they remain unchanged for years, even after the instrument has been taken apart. I call $b$ the inclination of the second axis to the plane of the circle, taken as positive when its left arm\footnote{as seen from behind the circle, which is assumed to be vertical.} is nearer to the circle. Finally I call $c$ the inclination of the second axis to a plane perpendicular to the first axis (which may be taken to be the plane of the azimuth circle), taken as positive when the left arm is higher above that plane.

\par\indent Let $M$ and $N$ be the intersections of the third axis with the second axis and with the line of sight, and let $O$ be a clearly visible object. The distance $MN$ has a definite value for each circle, to be found by measurement or otherwise. On our circle it is $\frac{1}{8}$ meter. When the line of sight is directed at the object $O$, I denote the angle $MON$ by $p$. Its exact formula is

\[ \sin p = \frac{MN}{MO}\cdot\cos a \]

\noindent for which one may always take $p = \dfrac{MN}{MO}\cdot 206265''$. It will not be difficult to determine $MO$ accurately enough to know $p$ to the nearest second. For objects at a considerable distance, $p$ becomes entirely imperceptible. That is no reason to prefer them to nearer objects in the following adjustment method if the latter offer better sighting points.

\par\indent I assume that in all the operations to be described here the circle itself is never turned about the third axis, so that the clamp on the drum is tightened firmly and its adjustment screw is left untouched. If desired, the clamping nut at the rear of the third axis may also be tightened a little more than usual.

\par\indent For brevity I call $I$ the angle between the plane through the second and third axes and the plane through the first and second axes. This angle thus determines how far the circle has turned about the second axis. The instrument is built to allow values of this angle from $0$ to $90^\circ$ and a little beyond both limits. I have already noted that the Reichenbach circles have no graduated scale for measuring this angle.

\par\indent The essential point of my method is that if the line of sight is aimed at an object with the instrument in any given position, the line of sight can be brought back to the object by combining a rotation about the second axis with a rotation of the alidade circle about the third. I assume that $I = i$ when the third axis is brought into one plane with the second axis and the straight line $MO$. The index of the alidade circle, on the other hand, is to read $k$ when the line of sight lies in one plane with the second and third axes, so that the line of sight and the second axis (left arm) have the same projection on the plane of the circle

\SourceLeaf{411}
\par\indent It is easy to see that if the line of sight is directed at the object when the index of the alidade circle reads $k-\varkappa$ and $I=i-\eta$, the line of sight will again be directed at the same object if $I=i+\eta$ is set and the index of the alidade circle is set to $k+\varkappa$. Of course, this second position is physically possible only if the second value of $I$ falls within the limits allowed by the construction of the instrument. It is also easy to show that the following equation holds exactly:

\[ \sin b \cos\varkappa = \cos b\, \mathrm{tang}(a+p) + \mathrm{cotang}\,\eta \sin\varkappa \]

\noindent The cosine of the angle between the second axis and the line $MO$

\[ = \sin b \sin(a+p) + \cos b \cos(a+p) \cos\varkappa \]

\noindent It follows readily that $\varkappa$ will either be very small or be very close to $180^\circ$. In the first case the angle in question will be very small, and in the second it will be close to $180^\circ$. That is, the object must lie nearly in the direction of the left arm of the second axis in the first case and nearly in the direction of the right arm in the second.

\par\indent We therefore have, to sufficient accuracy

\noindent in the first case \[ b - a - p = \frac{\varkappa}{\mathrm{tang}\,\eta} \]

\noindent in the second case \[ b + a + p = \frac{\varkappa - 180^\circ}{\mathrm{tang}\,\eta} \]

\par\indent If $\eta$ can be determined, two such trials give the values of $b-a-p$ and $b+a+p$, hence those of $b$ and $a+p$, or, insofar as $p$ can be regarded as known, that of $a$. At the same time each such trial obviously provides, without needing $\eta$, a determination of $k$. Since $\eta$ may be chosen freely here, it is best to make it as large as possible. Even so, it would be impossible to go much beyond $45^\circ$, and since the values of $I$ cannot be read on \textsc{Reichenbach} circles, one must restrict oneself to values for which the position can be obtained in some other way. I shall show below how to obtain with the greatest precision the position for which $I=0$, that is, where the first. second and third axes lie in one plane. The plane of the circle can also be made parallel to the first axis by setting both vertically in the usual way. These two positions can be marked with two lines drawn on the quadrant that carries the clamp for rotation about the second axis, along the edge of one of the supports carrying that axis, so that both positions can always be found again quickly and with sufficient precision. For brevity I shall call the first position \emph{horizontal} and the second \emph{vertical}. For the latter the exact relation is

\[ \cos I = -\,\mathrm{tang}\,b\,.\,\mathrm{tang}\,c \]

\noindent Thus $I$ will not differ appreciably from $90^\circ$, and since the difference is set $=2\eta$, the exact result will be

\[ \mathrm{tang}\,\eta = \BareRoot{\dfrac{1+\mathrm{tang}\,b\,.\,\mathrm{tang}\,c}{1-\mathrm{tang}\,b\,.\,\mathrm{tang}\,c}} = \BareRoot{\dfrac{\cos(b-c)}{\cos(b+c)}} \]

\noindent and so it will differ imperceptibly from $\eta$.

\SourceLeaf{412}
\par\indent It follows that if, in the first trial, with the object lying nearly in the direction of the left arm of the second axis, the index of the alidade circle read $h$ in the horizontal position of the circle and $\upsilon$ in the vertical position, while in the second trial, with the object lying nearly in the direction of the right arm of the second axis, the index of the alidade circle read $180^\circ+h'$, $180^\circ+\upsilon'$ respectively, we shall have:

\[ \begin{aligned} b-a-p &= \upsilon-k = k-h = \tfrac{1}{2}(\upsilon-h) \\ b+a+p &= \upsilon'-k = k-h' = \tfrac{1}{2}(\upsilon'-h') \end{aligned} \]

\noindent and consequently

\[ \begin{aligned} a &= \tfrac{1}{4}(\upsilon'-h'-\upsilon+h)-p \\ b &= \tfrac{1}{4}(\upsilon'-h'+\upsilon-h) \end{aligned} \]

\noindent At the same time we obtain the two values of $k$

\[ \tfrac{1}{2}(\upsilon+h) \quad\text{and}\quad \tfrac{1}{2}(\upsilon'+h') \]

\noindent whose close agreement will serve to confirm the correctness of the operations.

\par\indent My procedure thus runs as follows. The object I choose must present a sharply defined sighting point and be no more than a few degrees from the horizon. I set up the instrument so that one of its feet lies nearly in the same plane as the first axis and the object. Then, judging by eye, I bring the second axis into this plane as well by turning in azimuth, with its left arm facing the object at first. Next I set the circle horizontally and bring the line of sight onto the object by using the screw at the foot just mentioned and by moving the alidade. The circle is then set vertically, and the line of sight is brought back onto the object by turning in azimuth and moving the alidade. This completes the first rough setting. The circle is set horizontally again. If moving the alidade is insufficient to bring the line of sight onto the object, the foot screw is used to assist. The circle is then set vertically once more, and the line of sight is brought onto the object by moving the alidade, with the help of the azimuth screw if needed. These two operations are alternated until moving the alidade alone suffices for precise sighting in the horizontal position. The readings on the alidade circle give $h$ and $\upsilon$.

\par\indent Next, after the last vertical setting, the azimuth clamp is released and the circle is turned $180^\circ$ about the first axis, so that the right arm of the second axis now points nearly toward the object. The line of sight is brought precisely onto the object by moving the alidade (through nearly $180^\circ$) and making a fine adjustment with the azimuth screw. Here too I read the index of the alidade circle, which may read $180^\circ+V$. This is the beginning of the second series of operations. The circle is set horizontally, and the object is sighted by adjusting the foot screw and moving the alidade. At the new vertical setting the azimuth and the alidade are adjusted. The operations are again alternated until moving the alidade alone suffices to sight the object in both positions. The readings on the alidade circle here give $180^\circ+h'$ and $180^\circ+\upsilon'$. This completes the measurements.

\par\indent The reading $180^\circ+V$ mentioned above also serves to determine $c$. Evidently, when the circle is vertical, the line of sight will lie in the same plane as the first and third axes and will point upward when the index of the alidade circle reads $\tfrac{1}{2}(\upsilon+V)-90^\circ$. In other words, the projection of the line of sight on the plane of the circle will then be parallel to the first axis. The projection of the second axis (left arm) on that same plane therefore makes the angle $k+90^\circ-\tfrac{1}{2}(\upsilon+V)$ with the radius parallel to the first axis. Hence

\SourceLeaf{413}
\[ \sin c = \cos b\,\sin\bigl(\tfrac{1}{2}(\upsilon+V)-k\bigr) \]

\noindent for which

\[ c = \tfrac{1}{2}(\upsilon+V)-k = \tfrac{1}{2}(V-h) \]

\noindent may be taken.

\par\indent This procedure thus gives the error $a$ in seconds. To remove it when $a$ is positive, the crosshairs must be moved downward from the circle, or from the objective toward the eyepiece if a prismatic eyepiece is used. To visualize the amount of the shift, one can either estimate by eye how the shift compares with a suitable object of known diameter, or first aim the line of sight precisely at a fine target and then move the alidade circle by $a$. The crosshairs must then be shifted until the object appears on the diagonal. It goes without saying that this delicate operation must be done carefully and that any displacement of the circle itself must be avoided during it. In any event, repeating the method above provides a way to check whether the adjustment has worked.

\par\indent It remains for me to show how to set the circle in the horizontal position, where $I=0$. Insofar as $c$ and $b$ are small, the setting need not be all that exact and one could make do with judging by eye. The following procedure, however, achieves the setting with the utmost precision.

\par\indent Once $k$ has been determined with the drum fixed, first set the circle horizontally by eye, set the index of the alidade circle to $k+l$, where $l$ is an arbitrary angle, preferably $90^\circ$, and bring a clearly visible object onto the line of sight by turning in azimuth and adjusting a foot. Then set the index of the alidade circle to $k-l$. If the horizontal position happens to have been set exactly, a turn in azimuth alone will bring the line of sight precisely back onto the object. Otherwise the line of sight will pass above or below it. Then turn in azimuth to bring the object as close as possible to the line of sight (that is, onto the thread parallel to the third axis, or, with a prismatic eyepiece, onto the thread parallel to the telescope). Bring the object halfway to the line of sight by turning the circle about the second axis, and correct the other half with a foot screw, using the azimuth screw to help if needed, so that the object again lies exactly on the line of sight. To check the accuracy of the halving, set the index of the alidade circle to $k+l$ again and see whether a turn in azimuth alone now brings the object precisely back onto the line of sight. If it does not, again correct half the error, which will certainly now be much smaller, with the worm screw that makes the fine movement about the second axis. Continue until fully satisfied.

\par\indent If the first setting made by eye is too far off, it will be hard to make the first halving reasonably accurate, because the turn in azimuth will fail to bring the object into the field of view at all if $l$ has been taken large. To save much trial and error, it is therefore better to take $l$ small at first, say $=10^\circ$, and bring the circle into the desired position by the method just described. The setting can then be refined by repeating the procedure with a large value of $l$.

\par\indent If an instrument happens to have $b=c$, the plane of the circle in the horizontal position will be exactly perpendicular to the first axis. Once an object is on the line of sight, it can be brought back onto it by turning in azimuth alone after any change in the alidade setting. \SourceLeaf{414} But if $b$ and $c$ are unequal, it is easy to see that this is impossible. The plane of the circle cannot then be set exactly horizontal in the \emph{strict} sense unless the first axis is tilted a little. On our \textsc{Reichenbach} circle $c=+5'40''$, $b=+28''$, as the means of several trials.

\par\indent I have already noted that determining $k$ is quite independent of knowing particular values of $I$. One need only bring an object onto the line of sight twice by combining a movement of the alidade with a turn about the second axis, without touching the azimuth or drum screw, and take the mean of the two positions of the index of the alidade circle. It is best to arrange for as large a turn about the second axis as possible. For this purpose I first make the rough setting described above. Then, with the circle approximately horizontal and the object exactly on the line of sight, I turn the circle about the second axis until the arc described by the object in the field of view brings it back onto the thread parallel to the plane of the circle (with a prismatic eyepiece, onto the thread that represents it, that is, the one perpendicular to the telescope). I then move the alidade to bring the object exactly onto the cross. If the object \emph{remained} exactly on the cross while turning about the second axis, that would prove that the object lay on the second axis itself, that consequently $b=a+p$, and that the index position itself was $=k$.

\par\indent I would also like to mention here the procedure I find useful for setting the plane of the circle exactly vertical. The hanging spirit level is meant for this on \textsc{Reichenbach} circles. It can yield great accuracy, but using it is rather cumbersome, because the telescope must be removed, and it requires a very stable setup and a great deal of care. I use the mercury horizon and choose an object as high as possible (it need not be far away, since it does no harm in this operation if the eyepiece tube has to be drawn well out for a clear view). I set the circle vertically by eye, bring it into the vertical plane of the object, and aim exactly at it. The telescope is then lowered until the image of the object can be seen in the mercury horizon. If the line of sight falls exactly on it, the plane of the circle is vertical. Otherwise correct one half with either a foot screw or the worm screw, and the other half with the azimuth screw. If the halving has been exact, the line of sight of the telescope, when raised again, will meet the object itself exactly. Otherwise again correct one half of the deviation with the foot screw or worm screw and the other half with the azimuth screw. If the elevation of the object is not too small and it can be sighted with sufficient precision, this method too can yield very high accuracy, far more than is necessary unless one wants to measure \emph{very} small zenith distances. As is well known, however, it is inadvisable to use the \textsc{Borda} circle for zenith distances that are \emph{too small}.

